% Part: sets-functions-relations
% Chapter: functions
% Section: partial-functions

\documentclass[../../../include/open-logic-section]{subfiles}


\begin{document}

\olfileid[pt]{sfr}{fun}{par}

\olsection{Funções Parciais}

\begin{explain}
Às vezes é útil relaxar a definição de função de modo que não seja
exigido que a saída da função esteja definida para todas as
entradas possíveis. Tais mapeamentos são chamados de \emph{funções parciais}.
\end{explain}

\begin{defn}
Uma \emph{função parcial} $f \colon A \pto B$ é um mapeamento que
associa a cada !!{element} de~$A$ no máximo um !!{element} de~$B$.
Se $f$ associa um !!{element} de~$B$ a $x \in A$, dizemos que $f(x)$ está
\emph{definido}, e caso contrário \emph{indefinido}. Se $f(x)$ está definido,
escrevemos $f(x) \fdefined$, caso contrário $f(x) \fundefined$. O
\emph{domínio} de uma função parcial~$f$ é o subconjunto de~$A$ onde ela
está definida, isto é, $\dom{f} = \Setabs{x \in A}{f(x) \fdefined}$.
\end{defn}

\begin{ex}
Toda função $f\colon A \to B$ é também uma função parcial. Funções parciais
que estão definidas em todo~$A$---isto é, o que até agora
chamávamos simplesmente de função---também são chamadas de funções \emph{totais}.
\end{ex}

\begin{ex}
A função parcial $f \colon \Real \pto \Real$ dada por $f(x) = 1/x$
é indefinida para $x = 0$ e definida em todo o restante.
\end{ex}

\begin{prob}
Dada $f\colon A \pto B$, defina a função parcial $g\colon B \pto
A$ por: para qualquer $y \in B$, se existe um único $x \in A$ tal que
$f(x) = y$, então $g(y) = x$; caso contrário $g(y) \fundefined$.  Mostre que
se $f$ é injetiva, então $g(f(x)) = x$ para todo $x \in \dom{f}$ e
$f(g(y)) = y$ para todo $y \in \ran{f}$.
\end{prob}

\begin{defn}[Gráfico de uma função parcial]
Seja $f\colon A \pto B$ uma função parcial. O \emph{gráfico} de~$f$
é a relação $R_f \subseteq A \times B$ definida por
\[
R_f = \Setabs{\tuple{x,y}}{f(x) = y}.
\]
\end{defn}

\begin{prop}
Suponha que $R \subseteq A \times B$ tem a propriedade de que sempre que $Rxy$
e $Rxy'$ então $y = y'$.  Então $R$ é o gráfico da função parcial
$f\colon A \pto B$ definida por: se existe um $y$ tal que
$Rxy$, então $f(x) = y$, caso contrário $f(x) \fundefined$.  Se $R$ também é
\emph{serial}, isto é, para cada $x \in A$ existe um $y \in B$ tal que
$Rxy$, então $f$ é total.
\end{prop}

\begin{proof}
Suponha que existe um $y$ tal que $Rxy$.  Se existisse outro $y'
\neq y$ tal que $Rxy'$, a condição sobre $R$ seria
violada. Portanto, se existe um $y$ tal que $Rxy$, esse $y$ é
único, e assim $f$ é bem definida.  Obviamente, $R_f = R$ e $f$ é
total se~$R$ é serial.
\end{proof}

\end{document}
